\chapter{Reverse Architecture Documentation \\
\small{\textit{-- System Architecture and Component Analysis}}
\index{Architecture}
\index{Chapter!Architecture}
\label{Chapter::Architecture}}

This chapter presents the reverse-engineered architectural documentation of the DevOps Final Project, extracted automatically from source code and CI/CD pipeline configuration.

\section{Architecture Overview}
\label{sec:ArchOverview}

The DevOps Final Project is a hybrid system consisting of four main components:

\begin{enumerate}
    \item \textbf{Frontend Web Application} - HTML, CSS, and JavaScript files
    \item \textbf{Documentation System} - LaTeX-based technical documentation
    \item \textbf{Build Automation} - Compilation and metadata generation
    \item \textbf{CI/CD Pipeline} - GitHub Actions workflow orchestration
\end{enumerate}

\section{Frontend Layer}
\label{sec:FrontendLayer}

The frontend layer consists of three main files that form a web application:

\begin{table}[h]
\centering
\begin{tabular}{|l|l|l|}
\hline
\textbf{File} & \textbf{Type} & \textbf{Purpose} \\
\hline
index.html & HTML5 & Main web interface and document structure \\
\hline
style.css & CSS3 & Visual styling and responsive layout \\
\hline
demo.js & JavaScript & Interactive functionality and DOM manipulation \\
\hline
\end{tabular}
\caption{Frontend Components}
\label{table:frontend-components}
\end{table}

\subsubsection{Dependencies}

\begin{itemize}
    \item \texttt{index.html} depends on \texttt{style.css} for styling
    \item \texttt{index.html} depends on \texttt{demo.js} for interaction
    \item \texttt{demo.js} manipulates the DOM defined in \texttt{index.html}
\end{itemize}

\section{Documentation Layer}
\label{sec:DocumentationLayer}

The documentation layer is built on LaTeX and the Cornell thesis style:

\subsection{Core Documents}

\begin{table}[h]
\centering
\begin{tabular}{|l|l|l|}
\hline
\textbf{File} & \textbf{Type} & \textbf{Purpose} \\
\hline
itManual.tex & LaTeX & Master document that includes all chapters \\
\hline
prologue.tex & LaTeX & Title page, abstract, acknowledgments \\
\hline
itProjectDemo.tex & LaTeX & DevOps metrics and build information \\
\hline
itArchitecture.tex & LaTeX & System architecture and component analysis \\
\hline
\end{tabular}
\caption{Main Documentation Files}
\label{table:doc-files}
\end{table}

\subsection{Supporting Files}

The documentation system includes 17+ chapter files covering various aspects of the DevOps infrastructure including:

\begin{itemize}
    \item \texttt{itIntroduction.tex} - Project introduction
    \item \texttt{itAWSDeployment.tex} - AWS infrastructure setup
    \item \texttt{itGithub.tex} - GitHub configuration and workflows
    \item \texttt{itJenkinsWithPytest.tex} - Jenkins CI/CD setup
    \item \texttt{itPrometheusGrafana.tex} - Monitoring and logging systems
    \item \texttt{itLoadBalancer.tex} - Load balancing configuration
    \item \texttt{itKanbanSetup.tex} - Project management setup
    \item And more...
\end{itemize}

\section{Build Automation}
\label{sec:BuildAutomation}

The build automation layer handles LaTeX compilation and metadata generation:

\subsection{Version Generation}

Version information is automatically extracted from the git repository:

\begin{table}[h]
\centering
\begin{tabular}{|l|l|}
\hline
\textbf{Information} & \textbf{Source} \\
\hline
Version & Git tags or commit count \\
\hline
Commit Hash & Latest git commit SHA-1 \\
\hline
Author & Committer name from git \\
\hline
Timestamp & Commit date and time \\
\hline
Branch & Current git branch \\
\hline
Total Commits & Number of commits in repository \\
\hline
\end{tabular}
\caption{Automatically Captured Build Information}
\label{table:build-info}
\end{table}

\subsection{LaTeX Compilation}

The compilation process runs three passes to ensure proper cross-references:

\begin{enumerate}
    \item \textbf{First Pass}: Generate references and build reference list
    \item \textbf{Second Pass}: Build table of contents and index
    \item \textbf{Third Pass}: Final compilation with all references resolved
\end{enumerate}

\section{CI/CD Pipeline Architecture}
\label{sec:CIPipeline}

The CI/CD pipeline is orchestrated by GitHub Actions with multiple jobs:

\begin{table}[h]
\centering
\begin{tabular}{|l|l|l|}
\hline
\textbf{Workflow} & \textbf{Trigger} & \textbf{Purpose} \\
\hline
ci-cd.yml & Push to main/develop & Build, test, package, and release \\
\hline
\end{tabular}
\caption{GitHub Actions Workflows}
\label{table:workflows}
\end{table}

\subsection{Workflow Jobs}

The CI/CD pipeline includes the following main jobs:

\begin{enumerate}
    \item \textbf{Build Job} - Compiles source code, generates artifacts
    \item \textbf{Build LaTeX Job} - Runs LaTeX compilation with version info
    \item \textbf{Test Job} - Runs automated tests (validation)
    \item \textbf{Package and Release Job} - Creates deployment packages
\end{enumerate}

\subsection{Architecture Generation}

On each build, the system automatically generates:

\begin{itemize}
    \item \textbf{PlantUML Diagrams} - Component architecture specifications
    \item \textbf{Graphviz Diagrams} - Visual architecture representations
    \item \textbf{Source Analysis} - Code structure documentation
    \item \textbf{Component Maps} - System component relationships
\end{itemize}

\section{Data Flow}
\label{sec:DataFlow}

The data flows through the system as follows:

\begin{enumerate}
    \item Developer commits code and pushes to GitHub
    \item GitHub detects push and triggers workflows
    \item Build system extracts metadata from git
    \item LaTeX compiler generates diagrams and documentation
    \item PDF is created with all build information
    \item Artifacts are uploaded and released on GitHub
\end{enumerate}

\section{Build Output Artifacts}
\label{sec:OutputArtifacts}

The build system generates and stores the following artifacts:

\begin{table}[h]
\centering
\begin{tabular}{|l|l|l|}
\hline
\textbf{Artifact} & \textbf{Format} & \textbf{Storage} \\
\hline
itManual.pdf & PDF & GitHub Artifacts (30 days) \\
\hline
version.txt & Text & Auto-committed to repo \\
\hline
buildinfo.tex & LaTeX & Auto-committed to repo \\
\hline
buildinfo.json & JSON & GitHub Artifacts \\
\hline
architecture.png & PNG & Generated during build \\
\hline
architecture.pdf & PDF & Generated during build \\
\hline
devops-final-project-[ver].zip & ZIP & GitHub Releases \\
\hline
\end{tabular}
\caption{Build Output Artifacts}
\label{table:output-artifacts}
\end{table}

\section{Architecture Patterns}
\label{sec:ArchPatterns}

\subsection{Separation of Concerns}

The architecture maintains clear separation between:

\begin{itemize}
    \item \textbf{Frontend} - Web interface and user interaction
    \item \textbf{Documentation} - Technical and process documentation
    \item \textbf{Automation} - Build and compilation scripts
    \item \textbf{CI/CD} - Orchestration and release management
\end{itemize}

\subsection{Automation and Reproducibility}

\begin{itemize}
    \item All builds are automated via GitHub Actions
    \item Version information is automatically extracted from git
    \item Documentation is regenerated on every commit
    \item Artifacts are automatically stored and released
    \item Diagrams are generated from source code analysis
\end{itemize}

\subsection{Metadata Injection}

Build information is dynamically injected into:

\begin{itemize}
    \item PDF footer (version number on every page)
    \item Project Demo chapter (full build details)
    \item JSON metadata files (programmatic access)
    \item LaTeX commands (document integration)
    \item Architecture diagrams (system visualization)
\end{itemize}

\section{Technology Stack}
\label{sec:TechStack}

\begin{table}[h]
\centering
\begin{tabular}{|l|l|l|}
\hline
\textbf{Component} & \textbf{Technologies} & \textbf{Purpose} \\
\hline
Frontend & HTML5, CSS3, JavaScript & Web interface and interaction \\
\hline
Documentation & LaTeX, Cornell.cls, BibTeX & Technical documentation \\
\hline
Build Tools & pdflatex, sed, bash & Compilation and scripting \\
\hline
Diagrams & Graphviz, PlantUML & Architecture visualization \\
\hline
CI/CD & GitHub Actions, YAML & Workflow orchestration \\
\hline
Version Control & Git, GitHub & Source and documentation control \\
\hline
Packaging & ZIP, tar & Artifact distribution \\
\hline
\end{tabular}
\caption{Technology Stack}
\label{table:tech-stack}
\end{table}

\section{Deployment Architecture}
\label{sec:DeploymentArch}

The system uses a multi-stage deployment approach:

\begin{table}[h]
\centering
\begin{tabular}{|l|l|}
\hline
\textbf{Environment} & \textbf{Description} \\
\hline
Development & Local development on developer machine \\
\hline
Repository & GitHub repository with version control \\
\hline
CI/CD Server & GitHub Actions for build and test \\
\hline
Artifact Storage & GitHub Actions artifact storage (30 days) \\
\hline
Release & GitHub Releases for versioned releases \\
\hline
\end{tabular}
\caption{Deployment Architecture}
\label{table:deployment-arch}
\end{table}

\section{System Component Relationships}
\label{sec:ComponentRelations}

The key relationships between system components:

\begin{itemize}
    \item \textbf{HTML depends on CSS and JS} - Frontend layer dependencies
    \item \textbf{LaTeX depends on buildinfo} - Documentation includes metadata
    \item \textbf{PDF depends on LaTeX} - Compilation produces output
    \item \textbf{Workflows orchestrate builds} - CI/CD controls processes
    \item \textbf{Git metadata feeds diagrams} - Source code drives documentation
\end{itemize}

\section{Automatic Documentation Synchronization}
\label{sec:AutoSync}

The system maintains documentation synchronization through:

\begin{enumerate}
    \item \textbf{Git Hooks} - Trigger builds on push events
    \item \textbf{Automated Metadata Extraction} - Version info from git
    \item \textbf{Dynamic Diagram Generation} - Extracted from source code
    \item \textbf{PDF Regeneration} - On every commit
    \item \textbf{Artifact Management} - Automatic storage and retention
\end{enumerate}

This ensures that documentation always reflects the current state of the codebase and deployment pipeline.

\section{Key Features and Benefits}
\label{sec:Features}

\begin{itemize}
    \item \textbf{Automated Documentation} - Docs generated on every commit
    \item \textbf{Version Tracking} - Semantic versioning from git
    \item \textbf{Build Automation} - Fully automated LaTeX compilation
    \item \textbf{CI/CD Integration} - GitHub Actions orchestration
    \item \textbf{Reverse Documentation} - Architecture extracted from source
    \item \textbf{Artifact Management} - Automatic storage and release
    \item \textbf{Metadata Injection} - Build info embedded in PDF
    \item \textbf{Diagram Generation} - Visual architecture documentation
\end{itemize}

\section{Conclusions}
\label{sec:ArchConclusions}

The DevOps Final Project demonstrates a mature, automated approach to software documentation:

\begin{itemize}
    \item Documentation is treated as code and version-controlled
    \item Build information is automatically captured and embedded
    \item Architecture is visualized through automated diagram generation
    \item All processes are reproducible and auditable
    \item System maintains self-documenting properties
\end{itemize}

This architecture enables continuous documentation generation, ensuring that technical documentation stays synchronized with source code and build processes at all times.

\noindent
\textit{This chapter was automatically included from the system architecture documentation.}

